\chapter{Analisi del problema}

\section{L'algoritmo}
Le hash table nascono per rispondere all'esigenza di memorizzare grandi quantità di dati garantendo, nel caso medio, una complessità 
temporale pari a $O(1)$ per le operazioni di \texttt{search}, \texttt{insert} e \texttt{delete}.

L'idea di base consiste nel memorizzare dati organizzati in coppie chiave-valore all'interno di un array. Tale struttura consente l'accesso 
diretto agli elementi in tempo costante, sfruttando un indice calcolato tramite una funzione apposita, detta funzione di hash.

La funzione di hash ha il compito di trasformare una chiave in un indice dell'array in tempo $O(1)$. Una volta ottenuto l'indice, 
è possibile accedere direttamente alla posizione corrispondente nell'array, mantenendo quindi il tempo di accesso costante.

Affinché tutte le operazioni possano essere eseguite in tempo $O(1)$ anche nel caso peggiore, sarebbe necessario disporre di una funzione 
di hash perfetta, in grado di evitare completamente le collisioni (ovvero i casi in cui chiavi diverse producono lo stesso indice), 
e di uno spazio di memoria illimitato, così da non dover mai effettuare operazioni di ridimensionamento della struttura.

Tuttavia, tali condizioni ideali non sono realisticamente raggiungibili. Nella pratica, è necessario adottare opportune tecniche di gestione 
delle collisioni e strategie di ridimensionamento (resize). Grazie a queste soluzioni, è possibile garantire che, nel caso medio, 
le operazioni fondamentali mantengano una complessità $O(1)$, mentre nel caso peggiore possano degradare fino a $O(n)$.

\section{Strategie di gestione conflitti}
Nel lavoro effettuato sono state prese in considerazione quattro strategie di gestione dei conflitti, ovvero
\texttt{OPEN\_ADDRESSING\_LINEAR\_PROBING}, \break
\texttt{OPEN\_ADDRESSING\_DOUBLE\_HASHING},
\texttt{CHAINING} e \texttt{CUCKOO}.

\subsection{Open Addressing: Linear Probing}

L’open addressing è una tecnica di risoluzione delle collisioni in cui ogni elemento viene memorizzato direttamente nell’array dei bucket 
e la gestione delle collisioni avviene tramite probing (sondaggio).

Il \textit{linear probing} consiste nell’esaminare i bucket in modo sequenziale, partendo dalla posizione calcolata dalla funzione di hash e 
proseguendo linearmente (posizioni successive), finché non si trova la posizione desiderata.

Quando si deve inserire un nuovo elemento, i bucket vengono esaminati a partire dalla posizione calcolata dalla funzione di hash e si procede linearmente fino a trovare uno slot libero.

Quando si cerca un elemento, si segue la stessa sequenza di probing, finché non si trova il record desiderato oppure uno slot vuoto, 
il che indica che l’elemento non è presente nella tabella.

Questa tecnica ha complessità nel caso medio pari a $O(1)$, mentre nel caso peggiore è $O(n)$, in quanto potrebbe essere necessario scandire 
l’intero array. Le prestazioni dipendono fortemente dal fattore di carico $\alpha = \frac{n}{m}$ della tabella, che rappresenta il rapporto 
tra il numero di celle occupate e il numero di posizioni disponibili. Tipicamente, tale valore viene mantenuto nell’intervallo 0.6–0.75 
per garantire buone prestazioni.

\subsection{Open Addressing: Double Hashing}

L’open addressing è una tecnica di risoluzione delle collisioni in cui ogni elemento viene memorizzato direttamente nell’array dei bucket 
e la gestione delle collisioni avviene tramite probing (sondaggio).

Il \textit{double hashing} è una tecnica di probing avanzata, in cui la sequenza di posizioni da esaminare dipende da due funzioni di hash. 
La prima funzione di hash $h_1(k)$ determina la posizione iniziale del bucket, mentre la seconda funzione $h_2(k)$ definisce l’incremento 
per gli step successivi. Le posizioni da esaminare vengono calcolate come:
\[
(h_1(k) + i \cdot h_2(k)) \mod m, \quad i = 0, 1, 2, \dots
\]

Quando si deve inserire un nuovo elemento, si parte dalla posizione calcolata da $h_1(k)$ e, in caso di collisione, si procede seguendo 
la sequenza definita dalla combinazione di $h_1$ e $h_2$ fino a trovare uno slot libero.

Quando si cerca un elemento, si segue la stessa sequenza di probing. La ricerca termina quando si trova il record desiderato oppure uno 
slot vuoto, indicando che l’elemento non è presente nella tabella.

È fondamentale che $h_2(k)$ sia coprima con la dimensione dell’array $m$, altrimenti non tutte le posizioni disponibili potrebbero essere 
visitabili. Se $h_2(k)$ e $m$ non sono coprime, alcune posizioni rimarrebbero escluse dalla sequenza di probing, riducendo l’efficienza 
della tabella e aumentando il rischio di collisioni non risolvibili.

Il \textit{double hashing} riduce significativamente il problema del \textit{primary clustering} presente nel linear probing, poiché la 
sequenza di probing varia a seconda della chiave.

Questa tecnica ha complessità nel caso medio pari a $O(1)$, mentre nel caso peggiore può arrivare a $O(n)$. Le prestazioni dipendono fortemente 
dal fattore di carico $\alpha = \frac{n}{m}$, che rappresenta il rapporto tra il numero di celle occupate e il numero di posizioni disponibili. 
Tipicamente, tale valore viene mantenuto nell’intervallo 0.6–0.75 per garantire buone prestazioni.

\subsection{Chaining}

Il \textit{chaining} è una tecnica di risoluzione delle collisioni in cui ogni bucket dell’array non contiene un solo elemento, 
ma una lista (tipicamente una linked list) di tutti gli elementi che hashano nella stessa posizione.

Quando si deve inserire un nuovo elemento, si calcola la posizione tramite la funzione di hash $h(k)$ e si aggiunge l’elemento 
alla lista presente in quel bucket. In questo modo, anche se più elementi hashano nella stessa posizione, tutti possono essere memorizzati 
senza dover spostare altri elementi nell’array.

Quando si cerca un elemento, si calcola la posizione tramite $h(k)$ e si scorre la lista associata al bucket fino a trovare il record 
desiderato o fino a raggiungere la fine della lista, che indica che l’elemento non è presente.

Il chaining evita il problema del \textit{primary clustering} presente nelle tecniche di open addressing, poiché ogni bucket può 
contenere un numero illimitato di elementi hashati nella stessa posizione. Tuttavia, le prestazioni dipendono dalla lunghezza media delle liste.

La complessità nel caso medio per ricerca, inserimento o cancellazione è $O(1 + \alpha)$, dove $\alpha = n/m$ è il fattore di carico della tabella, 
cioè il rapporto tra il numero totale di elementi $n$ e il numero di bucket $m$. Nel caso peggiore, se tutti gli elementi hashano nello 
stesso bucket, la complessità può arrivare a $O(n)$, ma questo scenario è raro se la funzione di hash distribuisce bene gli elementi.

Questa tecnica è particolarmente efficace quando il numero di elementi da memorizzare non è noto a priori o quando il fattore di carico 
può superare 1, cosa che non è possibile con l’open addressing senza ridimensionare la tabella. Tipicamente, tale valore viene mantenuto 
nell’intervallo 1–3 per garantire buone prestazioni.

\subsection{Cuckoo Hashing}

Il \textit{Cuckoo Hashing} è una tecnica di risoluzione delle collisioni che utilizza due (o più) funzioni di hash e due tabelle separate. 
Ogni elemento può essere memorizzato in una sola posizione per ciascuna funzione di hash, e in caso di collisione un elemento “scaccia” quello 
già presente nella posizione, il quale viene reinserito nella sua posizione alternativa calcolata dalla seconda funzione di hash.

Quando si deve inserire un nuovo elemento, si calcola la prima posizione tramite $h_1(k)$.  
- Se lo slot è libero, l’elemento viene inserito.  
- Se lo slot è occupato, l’elemento corrente viene “espulso” e inserito nella sua posizione alternativa calcolata da $h_2(k)$.  
Questo processo di spostamento può ripetersi fino a che tutti gli elementi trovano una posizione libera o fino a raggiungere un limite massimo di spostamenti.  
In quest’ultimo caso, la tabella viene ridimensionata o le funzioni di hash vengono ricalcolate.

Quando si cerca un elemento, basta controllare le posizioni calcolate da $h_1(k)$ e $h_2(k)$; se l’elemento non è presente in nessuna delle due, 
significa che non è nella tabella. Questa caratteristica garantisce che la ricerca sia sempre $O(1)$ nel caso peggiore.

Il Cuckoo Hashing evita il problema del \textit{primary clustering} presente nel linear probing, e mantiene un fattore di carico relativamente alto 
senza compromettere le prestazioni di ricerca. Tuttavia, le operazioni di inserimento possono richiedere più spostamenti e, in casi rari, 
il ridimensionamento della tabella.

La complessità media di ricerca è $O(1)$, così come nel caso peggiore, mentre l’inserimento ha complessità media $O(1)$ ma nel caso peggiore 
può richiedere più operazioni di spostamento o un ridimensionamento della tabella. Tipicamente, il fattore di carico $\alpha = n/m$ viene mantenuto 
nell’intervallo 0.4–0.6 per garantire buone prestazioni e ridurre il numero di cicli di spostamento.


\section{Strategie di Resizing}

Nel lavoro effettuato sono state prese in considerazione tre strategie di \textit{resize}, ovvero:  
\texttt{DOUBLE}, \texttt{CONSTANT} e \texttt{DYNAMIC}.  
Queste strategie definiscono come la dimensione della tabella hash viene aumentata quando il fattore di carico supera un certo valore soglia.

\subsection{Double}

La strategia \texttt{DOUBLE} consiste nel raddoppiare la dimensione della tabella ogni volta che il fattore di carico supera una soglia prefissata.  
Quando avviene il \textit{resize}, tutti gli elementi vengono re-inseriti nella nuova tabella utilizzando la stessa funzione di hash.  
Questo approccio è semplice e garantisce che il fattore di carico torni ad essere basso dopo ogni ridimensionamento.  
Tuttavia, il costo dell’operazione di raddoppio può essere elevato, poiché tutti gli elementi devono essere reinseriti.

\subsection{Constant}

La strategia \texttt{CONSTANT} prevede di aumentare la dimensione della tabella di un valore fisso ogni volta che il fattore di carico supera la soglia.  
Questo approccio distribuisce il costo del \textit{resize} in modo più uniforme rispetto al raddoppio, ma può comportare più operazioni 
di resize se la crescita del numero di elementi è elevata.  
È utile quando si prevede una crescita lineare del numero di elementi nella tabella.

\subsection{Dynamic}

La strategia \texttt{DYNAMIC} prevede il raddoppio della dimensione della tabella, come nella strategia \texttt{DOUBLE}, 
ma lo spostamento degli elementi avviene in modo incrementale.  
In pratica, non tutti gli elementi vengono copiati nella nuova tabella immediatamente, ma viene spostato un \textit{chunk} di elementi 
ad ogni operazione di inserimento o ricerca, fino a completare il trasferimento di tutti gli elementi dalla vecchia tabella a quella nuova.

Quando si cerca un elemento durante il periodo di trasferimento, si procede nel seguente modo:
\begin{itemize}
    \item Si cerca prima nella nuova tabella.
    \item Se non è presente, si cerca nella vecchia tabella.
    \item Se l’elemento è trovato nella vecchia tabella, viene copiato nella nuova.
\end{itemize}

Questa tecnica permette di distribuire il costo del ridimensionamento su più operazioni, evitando picchi di tempo elevati durante un singolo inserimento.  
Mantiene le prestazioni di ricerca e inserimento prossime a $O(1)$ anche durante il processo di resizing, 
rendendola più efficiente rispetto al raddoppio immediato della strategia \texttt{DOUBLE}.